\chapter{Complexity — ARM Sim}
%%% generated by the blueprint pipeline from TCSlib.Complexity.SpaceComplexity.Machines.ARMSim
\section{Overview}

Declarations of the Lean module \texttt{TCSlib.Complexity.SpaceComplexity.Machines.ARMSim}.

\section{Declarations}

\begin{definition}[Intermediate machine configurations]
\lean{Complexity.LogProg.Mid}
\label{Complexity.LogProg.Mid}
\leanok
\uses{Complexity.LogProg.Ph, Complexity.LogProg.RProg, Turing.Cfg}
Let $P$ be a program with $m$ work tapes and control states of the form
$(\lambda,\phi)$, let $W$ be an integer, and let $c$ be a configuration on input $x$.
The configuration is intermediate if every state currently represented by $c$ is not a
call state of $P$, and every work-tape head position lies between $-1$ and $W$,
inclusive.
\end{definition}

\begin{lemma}[Transition function of a compiled register program]
\lean{Complexity.LogProg.armProg_tr}
\label{Complexity.LogProg.armProg_tr}
\leanok
\uses{Complexity.LogProg.ARM, Complexity.LogProg.Ph, Complexity.LogProg.armProg, Complexity.LogProg.insTr}
\difficulty{1}
Let $A$ assign an instruction to each label in $\Lambda$, and let $l_0,l\in\Lambda$,
phase $ph$, input symbol $a\in\mathrm{Option}(\mathrm{Bool})$, and work-tape symbols
$w:\mathrm{Fin}(m)\to\mathrm{Option}(\mathrm{Bool})$ be given. The transition of the
compiled register program from initial label $l_0$, at state $(l,ph)$ with scanned
symbols $a$ and $w$, is the action prescribed by the instruction at $l$ for label $l$,
phase $ph$, and those scanned symbols.
\end{lemma}

\begin{lemma}[Call behavior of a compiled register program]
\lean{Complexity.LogProg.armProg_call}
\label{Complexity.LogProg.armProg_call}
\leanok
\uses{Complexity.LogProg.ARM, Complexity.LogProg.Ph, Complexity.LogProg.armProg, Complexity.LogProg.insCall}
\difficulty{1}
Let $A$ assign an instruction to each label in $\Lambda$, and let $l_0,l\in\Lambda$ and
$ph$ be a phase. The call specification of the register program compiled from $A$ with
initial label $l_0$, evaluated at $(l,ph)$, equals the call specification produced
directly by applying the instruction at $l$ to phase $ph$.
\end{lemma}

\begin{lemma}[Register tape at the seam configuration]
\lean{Complexity.LogProg.aseam_eq_regCfg}
\label{Complexity.LogProg.aseam_eq_regCfg}
\leanok
\uses{Complexity.LogProg.aseam, Complexity.LogProg.regCfg, Turing.FinTM.bufferTape}
\difficulty{4}
For any control state $l$, register assignment $v$, output word $out$, and register $r$,
the seam configuration encoding $l$, $v$, and $out$ is unchanged by replacing its $r$th
work tape with the tape storing the binary representation of $v(r)$ from cell zero and
setting that tape's head position to zero.
\end{lemma}

\begin{lemma}[Updating a register in a seam configuration]
\lean{Complexity.LogProg.regCfg_aseam}
\label{Complexity.LogProg.regCfg_aseam}
\leanok
\uses{Complexity.LogProg.aseam, Complexity.LogProg.regCfg, Turing.FinTM.bufferTape}
\difficulty{4}
Let $l,l'$ be states, let $v$ assign a natural number to each of $m$ registers, let
$\mathrm{out}$ be a Boolean word, let $r$ be a register, and let $n$ be a natural
number. Replacing register $r$'s tape in the seam configuration determined by
$l,v,\mathrm{out}$ with the binary word for $n$, setting its head to $0$, and changing
the state to $(l',\mathrm{start})$ gives the seam configuration determined by $l'$, the
function obtained from $v$ by setting $v(r)=n$, and the same output word.
\end{lemma}

\begin{lemma}[Register configuration satisfies the middle invariant]
\lean{Complexity.LogProg.mid_regCfg}
\label{Complexity.LogProg.mid_regCfg}
\leanok
\uses{Complexity.LogProg.ARM, Complexity.LogProg.Mid, Complexity.LogProg.Ph, Complexity.LogProg.armProg, Complexity.LogProg.aseam, Complexity.LogProg.regCfg}
\difficulty{4}
Let $m,d$ be nonnegative integers, let $\Lambda$ be a state set, and fix a Boolean input
$x$, an register machine $A$ with $m$ registers and $d$ deciders, an initial label
$l_0$, a label $l$, register values $v$, an output word $out$, a register $r$, a
phase-augmented state $s$, a tape content function $f:\mathbb
Z\to\mathrm{Option}(\{0,1\})$, and integers $q,W$. If the program associated to $A$ has
no decider call at state $s$, $-1\le q\le W$, and $0\le W$, then the configuration
obtained from the initial register configuration by setting its state to $s$, replacing
register $r$'s tape by $f$ with head at $q$, satisfies the middle invariant: it has no
decider call at its current state, and every register head lies between $-1$ and $W$,
inclusive.
\end{lemma}

\begin{lemma}[Mid-configuration bounds after an update]
\lean{Complexity.LogProg.mid_xCfg}
\label{Complexity.LogProg.mid_xCfg}
\leanok
\uses{Complexity.LogProg.ARM, Complexity.LogProg.Mid, Complexity.LogProg.Ph, Complexity.LogProg.armProg, Complexity.LogProg.aseam, Complexity.LogProg.xCfg}
\difficulty{4}
Let $A$ be a register program with $m$ registers, let $l_0,l$ be program states, let $v$
assign a natural number to each register, and let $out$ be a Boolean output word. Choose
a register $r$, a state $s$ paired with a phase, an input-head position $ip$, and
integers $q,W$. If the compiled program has no call at $s$, if $-1\le q\le W$, and if
$0\le W$, then the configuration obtained by starting from the seam configuration for
$l,v,out$ and setting its state to $s$, input-head position to $ip$, and register-head
position at $r$ to $q$ is a mid-configuration: it has no call at its active state, and
every register-head position lies between $-1$ and $W$, inclusive.
\end{lemma}

\begin{lemma}[Resetting a configuration at the initial seam]
\lean{Complexity.LogProg.xCfg_aseam}
\label{Complexity.LogProg.xCfg_aseam}
\leanok
\uses{Complexity.LogProg.aseam, Complexity.LogProg.xCfg}
\difficulty{4}
Let $x$ be an input string, let $l,l'$ be states, let $v$ assign a natural number to
each of the $m$ registers, let $out$ be a Boolean output word, and let $r$ be a
register. Applying the configuration update that sets the state to
$(l',\mathrm{start})$, places the input head at position $1$, and moves register $r$'s
work-tape head to $0$ to the initial seam configuration with state $l$ yields the
initial seam configuration with state $l'$, with all other contents and positions as
specified by that seam configuration.
\end{lemma}

\begin{lemma}[Initial Configuration Is Preserved]
\lean{Complexity.LogProg.xCfg_aseam_self}
\label{Complexity.LogProg.xCfg_aseam_self}
\leanok
\uses{Complexity.LogProg.aseam, Complexity.LogProg.xCfg, Complexity.LogProg.xCfg_aseam}
\difficulty{1}
Let $m$ be a natural number, let $x$ be a finite bit string, let $\Lambda$ be a type,
and let $l\in\Lambda$, $v:\{0,\ldots,m-1\}\to\mathbb{N}$, and $\mathrm{out}$ be a finite
bit string. For every work-tape index $r<m$, the configuration initialized from $l$,
$v$, and $\mathrm{out}$ is unchanged when its state and input-head position are reset to
their initial values and work-tape head $r$ is set to position $0$.
\end{lemma}

\begin{lemma}[Head positions in a seam configuration]
\lean{Complexity.LogProg.eqCfg_aseam}
\label{Complexity.LogProg.eqCfg_aseam}
\leanok
\uses{Complexity.LogProg.aseam, Complexity.LogProg.eqCfg}
\difficulty{4}
For any input, states $l,l'$, register valuation $v$, output word, and registers $r,s$,
changing a seam configuration from state $(l,\mathrm{start})$ to $(l',\mathrm{start})$
by setting both register-head positions to zero gives exactly the seam configuration
with state $(l',\mathrm{start})$.
\end{lemma}

\begin{definition}[Simulation with intermediate configurations]
\lean{Complexity.LogProg.SimTo}
\label{Complexity.LogProg.SimTo}
\leanok
\uses{Complexity.LogProg.Mid, Complexity.LogProg.Ph, Complexity.LogProg.RProg, Complexity.LogProg.rrun, Turing.Cfg}
Let $P$ be a program with $m$ work tapes and control states of the form
$(\lambda,\phi)$, let an oracle assign Boolean answers to deciders on bit strings, let
$W$ be an integer, and let $c,c'$ be configurations on input $x$. The configuration $c'$
is reached from $c$ after some finite number $T$ of steps of the run of $P$ with that
oracle, and every configuration at a time $t<T$ is intermediate: no represented state is
a call state, and all work-tape head positions lie between $-1$ and $W$, inclusive.
\end{definition}

\begin{definition}[Bounded halting with output]
\lean{Complexity.LogProg.HaltsWith}
\label{Complexity.LogProg.HaltsWith}
\leanok
\uses{Complexity.LogProg.Mid, Complexity.LogProg.Ph, Complexity.LogProg.RProg, Complexity.LogProg.rrun, Turing.Cfg}
Let $P$ be a program with $m$ work tapes and states in $\Lambda\times\mathrm{Ph}$, let
an oracle assign a Boolean answer to each of $d$ deciders and Boolean-word queries, and
let $c$ be a configuration on input $x$ with output word $v$. The machine must halt
after some number $T$ of steps with output $v$ followed by the single symbol $b$; at
that time every work-tape head position must lie between $-1$ and $W$, inclusive, and
every configuration at a time $t<T$ must satisfy the condition
$\mathrm{Mid}(P,W,\cdot)$.
\end{definition}

\begin{lemma}[Simulation of Register Increment]
\lean{Complexity.LogProg.sim_inc}
\label{Complexity.LogProg.sim_inc}
\leanok
\uses{Complexity.LogProg.SimTo, Complexity.LogProg.armProg, Complexity.LogProg.armProg_call, Complexity.LogProg.armProg_tr, Complexity.LogProg.aseam, Complexity.LogProg.aseam_eq_regCfg, Complexity.LogProg.inc_run, Complexity.LogProg.mid_regCfg, Complexity.LogProg.regCfg_aseam}
\difficulty{5}
Let $A$ be a register-program instruction map with initial label $l_0$, let the oracle
assign a Boolean answer to each decider and bit string, and let $x$ be an input bit
string. Suppose labels $l,l'$ and a register $r$ satisfy
$A(l)=\operatorname{inc}(r,l')$. For any register valuation $v$ and output word, if the
length of the binary expansion of $v(r)+1$ is at most the integer bound $W$, then the
compiled program simulates this increment from the configuration with label $l$ and
valuation $v$ to the configuration with label $l'$ and valuation obtained by replacing
$v(r)$ by $v(r)+1$, preserving the input and output; throughout the simulation, every
intermediate configuration before the endpoint has no pending oracle call and all
work-head positions lie between $-1$ and $W$, inclusive.
\end{lemma}

\begin{lemma}[Simulation of Register Decrement]
\lean{Complexity.LogProg.sim_dec}
\label{Complexity.LogProg.sim_dec}
\leanok
\uses{Complexity.LogProg.SimTo, Complexity.LogProg.armProg, Complexity.LogProg.armProg_call, Complexity.LogProg.armProg_tr, Complexity.LogProg.aseam, Complexity.LogProg.aseam_eq_regCfg, Complexity.LogProg.dec_run, Complexity.LogProg.mid_regCfg, Complexity.LogProg.regCfg_aseam}
\difficulty{4}
Let a program have finitely many natural-number registers, labels in a set $\Lambda$,
and an oracle for its calls. Suppose its instruction at label $l$ is to decrement
register $r$ and continue at label $l'$. For any Boolean input $x$, initial register
values $v$, output word, and integer bound $W$ with the length of the binary word for
$v(r)$ at most $W$, the compiled program can run from the configuration with label $l$,
registers encoded by their binary words, all work heads at cell $0$, and the given
output to the corresponding configuration with label $l'$, register $r$ set to
$\max(v(r)-1,0)$, and the other registers and output unchanged. During this finite run,
every earlier configuration has no pending oracle call, and every work head position
lies between $-1$ and $W$, inclusive.
\end{lemma}

\begin{lemma}[Simulation of Register Clearing]
\lean{Complexity.LogProg.sim_clr}
\label{Complexity.LogProg.sim_clr}
\leanok
\uses{Complexity.LogProg.SimTo, Complexity.LogProg.armProg, Complexity.LogProg.armProg_call, Complexity.LogProg.armProg_tr, Complexity.LogProg.aseam, Complexity.LogProg.aseam_eq_regCfg, Complexity.LogProg.clr_run, Complexity.LogProg.mid_regCfg, Complexity.LogProg.regCfg_aseam}
\difficulty{4}
Let a program have finitely many natural-number registers, labels in a set $\Lambda$,
and an oracle for its calls. Suppose its instruction at label $l$ is to clear register
$r$ and continue at label $l'$. For any Boolean input $x$, initial register values $v$,
output word, and integer bound $W$ with the length of the binary word for $v(r)$ at most
$W$, the compiled program can run from the configuration with label $l$, registers
encoded by their binary words, all work heads at cell $0$, and the given output to the
corresponding configuration with label $l'$, register $r$ set to $0$, and the other
registers and output unchanged. During this finite run, every earlier configuration has
no pending oracle call, and every work head position lies between $-1$ and $W$,
inclusive.
\end{lemma}

\begin{lemma}[Simulation of Register Halving]
\lean{Complexity.LogProg.sim_half}
\label{Complexity.LogProg.sim_half}
\leanok
\uses{Complexity.LogProg.SimTo, Complexity.LogProg.armProg, Complexity.LogProg.armProg_call, Complexity.LogProg.armProg_tr, Complexity.LogProg.aseam, Complexity.LogProg.aseam_eq_regCfg, Complexity.LogProg.half_run, Complexity.LogProg.mid_regCfg, Complexity.LogProg.regCfg_aseam}
\difficulty{4}
Let a program have finitely many natural-number registers, labels in a set $\Lambda$,
and an oracle for its calls. Suppose its instruction at label $l$ is to replace register
$r$ by its integer quotient upon division by $2$ and continue at label $l'$. For any
Boolean input $x$, initial register values $v$, output word, and integer bound $W$ with
the length of the binary word for $v(r)$ at most $W$, the compiled program can run from
the configuration with label $l$, registers encoded by their binary words, all work
heads at cell $0$, and the given output to the corresponding configuration with label
$l'$, register $r$ set to $\lfloor v(r)/2\rfloor$, and the other registers and output
unchanged. During this finite run, every earlier configuration has no pending oracle
call, and every work head position lies between $-1$ and $W$, inclusive.
\end{lemma}

\begin{lemma}[Simulation of a State Transition]
\lean{Complexity.LogProg.sim_go}
\label{Complexity.LogProg.sim_go}
\leanok
\uses{Complexity.LogProg.Ph, Complexity.LogProg.SimTo, Complexity.LogProg.armProg, Complexity.LogProg.aseam, Complexity.LogProg.goAct, Complexity.LogProg.rrun_one, Complexity.LogProg.rrun_zero, Complexity.LogProg.rstep_noncall, Turing.Action.apply, Turing.Cfg, Turing.Cfg.workTapeSymbols, Turing.FinTM.bufferTape, Turing.MultiTapeTM.step}
\difficulty{4}
Let a program have finitely many natural-number registers, labels in a set $\Lambda$,
and an oracle for its calls. Suppose that at label $l$ it has no pending oracle call,
and that for every possible input symbol the machine transition from label $l$, when
each work head scans cell $0$ containing the binary encoding of its register value,
moves no heads, writes nothing, emits no output, and changes the label to a specified
target label. For any Boolean input $x$, register values $v$, output word, and integer
$W\ge 0$, the compiled program can run from the configuration with label $l$, those
register values, all work heads at cell $0$, and the given output to the configuration
with the target label and the same registers and output. During this finite run, every
earlier configuration has no pending oracle call, and every work head position lies
between $-1$ and $W$, inclusive.
\end{lemma}

\begin{lemma}[Simulation of a zero test]
\lean{Complexity.LogProg.sim_jz}
\label{Complexity.LogProg.sim_jz}
\leanok
\uses{Complexity.LogProg.SimTo, Complexity.LogProg.armProg, Complexity.LogProg.armProg_call, Complexity.LogProg.armProg_tr, Complexity.LogProg.aseam, Complexity.LogProg.bits_head_odd, Complexity.LogProg.insTr, Complexity.LogProg.sim_go, Turing.FinTM.bufferTape}
\difficulty{3}
Let $m$ be the number of registers, let $\Lambda$ be a set of program states, and let
$A$ assign an instruction to each state. Fix an initial state $l_*$, an oracle on
Boolean words, an input word $x$, register values
$v\colon\{0,\ldots,m-1\}\to\mathbb{N}$, an output word, and an integer $W\ge 0$. Suppose
the instruction at $l$ tests whether register $r$ is zero, with successor states $l_1$
and $l_0'$. Then the compiled program, starting in state $l$ with those register values
and output, reaches the configuration with state $l_1$ if $v(r)=0$ and state $l_0'$
otherwise; throughout the run before that configuration, no call state is visited and
every register head position lies between $-1$ and $W$, inclusive.
\end{lemma}

\begin{lemma}[Simulation of an oddness test]
\lean{Complexity.LogProg.sim_jodd}
\label{Complexity.LogProg.sim_jodd}
\leanok
\uses{Complexity.LogProg.SimTo, Complexity.LogProg.armProg, Complexity.LogProg.armProg_call, Complexity.LogProg.armProg_tr, Complexity.LogProg.aseam, Complexity.LogProg.bits_head_odd, Complexity.LogProg.insTr, Complexity.LogProg.sim_go, Turing.FinTM.bufferTape}
\difficulty{4}
Let $m$ be the number of registers, let $\Lambda$ be a set of program states, and let
$A$ assign an instruction to each state. Fix an initial state $l_*$, an oracle on
Boolean words, an input word $x$, register values
$v\colon\{0,\ldots,m-1\}\to\mathbb{N}$, an output word, and an integer $W\ge 0$. Suppose
the instruction at $l$ tests whether register $r$ is odd, with successor states $l_1$
and $l_0'$. Then the compiled program, starting in state $l$ with those register values
and output, reaches the configuration with state $l_1$ if $v(r)$ is odd and state $l_0'$
otherwise; throughout the run before that configuration, no call state is visited and
every register head position lies between $-1$ and $W$, inclusive.
\end{lemma}

\begin{lemma}[Simulation of Register Equality Test]
\lean{Complexity.LogProg.sim_jeq}
\label{Complexity.LogProg.sim_jeq}
\leanok
\uses{Complexity.LogProg.SimTo, Complexity.LogProg.armProg, Complexity.LogProg.armProg_call, Complexity.LogProg.armProg_tr, Complexity.LogProg.aseam, Complexity.LogProg.eqCfg, Complexity.LogProg.eqCfg_aseam, Complexity.LogProg.eq_run}
\difficulty{4}
Let $A$ be a register-program instruction map with initial label $l_0$, let the oracle
assign a Boolean answer to each decider and bit string, and let $x$ be an input bit
string. Suppose labels $l,l_1,l_0'$, distinct registers $r,s$, and the instruction
identity $A(l)=\operatorname{jeq}(r,s,l_1,l_0')$ are given. For any register valuation
$v$ and output word, if the length of the binary expansion of $v(r)$ is at most the
integer bound $W$, then the compiled program simulates this instruction from the
configuration with label $l$ and valuation $v$ to the configuration with the same
valuation and output and with label $l_1$ when $v(r)=v(s)$, and label $l_0'$ otherwise;
at every intermediate configuration before the endpoint, no oracle call is pending and
each work-head position lies between $-1$ and $W$, inclusive.
\end{lemma}

\begin{lemma}[One-step effect of an oracle call]
\lean{Complexity.LogProg.sim_call}
\label{Complexity.LogProg.sim_call}
\leanok
\uses{Complexity.LogProg.CallSpec, Complexity.LogProg.Mode, Complexity.LogProg.Ph, Complexity.LogProg.armProg, Complexity.LogProg.armProg_call, Complexity.LogProg.aseam, Complexity.LogProg.callSegs, Complexity.LogProg.regWords, Complexity.LogProg.rrun, Complexity.LogProg.rrun_one, Complexity.LogProg.rstep, Complexity.LogProg.tapeWord_bufferTape, Complexity.LogProg.vword, Turing.Cfg, Turing.FinTM.bufferTape}
\difficulty{3}
Let $A$ be an instruction program on $m$ registers with control states in $\Lambda$, let
the oracle answer Boolean queries to $d$ deciders, and fix its initial state $l_0$ and a
Boolean input $x$. Suppose the instruction at state $l$ is a call to decider $j$ in mode
$\mu$ with argument-register list $R$, branching to $l_1$ on a positive answer and to
$l_0'$ on a negative answer. For any register assignment $v$ and output word $o$, one
step of the program from the configuration in state $l$, at the start phase, with input
head at position one and register tapes storing the binary expansions of $v$, yields the
corresponding start-phase configuration with the same registers, input-head position,
and output, and state $l_1$ if the oracle accepts the virtual input formed from $x$,
mode $\mu$, and the words in registers $R$, or state $l_0'$ otherwise.
\end{lemma}

\begin{lemma}[Halting behavior of a return instruction]
\lean{Complexity.LogProg.sim_ret}
\label{Complexity.LogProg.sim_ret}
\leanok
\uses{Complexity.LogProg.HaltsWith, Complexity.LogProg.armProg, Complexity.LogProg.armProg_call, Complexity.LogProg.armProg_tr, Complexity.LogProg.aseam, Complexity.LogProg.insTr, Complexity.LogProg.retAct, Complexity.LogProg.rrun, Complexity.LogProg.rrun_one, Complexity.LogProg.rrun_zero, Complexity.LogProg.rstep_noncall, Turing.Action.apply, Turing.FinTM.bufferTape, Turing.MultiTapeTM.step}
\difficulty{4}
Let $m$ be the number of registers, let $\Lambda$ be a set of program states, and let
$A$ assign an instruction to each state. Fix an initial state $l_*$, an oracle on
Boolean words, an input word $x$, register values
$v\colon\{0,\ldots,m-1\}\to\mathbb{N}$, an output word, and an integer $W\ge 0$. If the
instruction at $l$ returns the Boolean value $b$, then the compiled program, started in
state $l$ with those register values and output, reaches a halting configuration whose
output is the given output word followed by $b$, with every register head position
between $-1$ and $W$, inclusive. At every earlier step no call state is visited and all
register head positions lie in that same interval.
\end{lemma}

\begin{lemma}[Rejection implies bounded halting]
\lean{Complexity.LogProg.haltsWith_of_rejects}
\label{Complexity.LogProg.haltsWith_of_rejects}
\leanok
\uses{Complexity.LogProg.HaltsWith, Complexity.LogProg.Rejects, Complexity.LogProg.armProg, Complexity.LogProg.aseam, Complexity.LogProg.mid_xCfg}
\difficulty{4}
Let $A$ be a register program with initial label $l_0$, let an oracle answer queries to
its deciders, and let $x$ be a Boolean input. For any label $l$, register assignment
$v$, output word, register $r$, and integer $W\ge 0$, if the run of the compiled program
from the configuration initialized by $x$, $v$, and that output rejects at register $r$
with head position reset to $0$, then it halts with output bit $0$ and all work-tape
heads in the interval $[-1,W]$. Before halting, every configuration is between ordinary
machine steps rather than oracle calls, and all its work-tape heads lie in $[-1,W]$.
\end{lemma}

\begin{lemma}[Reachability implies bounded simulation]
\lean{Complexity.LogProg.simTo_of_reaches}
\label{Complexity.LogProg.simTo_of_reaches}
\leanok
\uses{Complexity.LogProg.Reaches, Complexity.LogProg.SimTo, Complexity.LogProg.armProg, Complexity.LogProg.aseam, Complexity.LogProg.mid_xCfg}
\difficulty{3}
Let $A$ be an instruction program on $m$ registers, with control states in $\Lambda$,
and let an oracle answer Boolean queries to $d$ deciders. Fix an initial state $l_0$,
Boolean input $x$, states $l,l'\in\Lambda$, register values $v:\{0,\ldots,m-1\}\to\bbn$,
output word $o$, register $r$, and integer $W\ge 0$. Form the initial and final
configurations with states $l$ and $l'$, respectively, phase at its start, input head at
position one, every register tape storing the binary expansion of its assigned value
with head at zero, and output $o$. If the program reaches the final configuration from
the initial one and, at every strictly earlier time, is at an ordinary (non-call) state
with register $r$'s head at position zero, then the run is a simulation to the final
configuration in which every strictly earlier configuration is an ordinary state and
every register head position lies between $-1$ and $W$, inclusive.
\end{lemma}

\begin{lemma}[Bounded reachability implies simulation]
\lean{Complexity.LogProg.simTo_of_reachesB}
\label{Complexity.LogProg.simTo_of_reachesB}
\leanok
\uses{Complexity.LogProg.ReachesB, Complexity.LogProg.SimTo, Complexity.LogProg.armProg, Complexity.LogProg.aseam, Complexity.LogProg.mid_xCfg}
\difficulty{3}
Let $A$ be an instruction program on $m$ registers, with control states in $\Lambda$,
and let an oracle answer Boolean queries to $d$ deciders. Fix an initial state $l_0$,
Boolean input $x$, states $l,l'\in\Lambda$, register values $v:\{0,\ldots,m-1\}\to\bbn$,
output word $o$, register $r$, and integers $L,W$ with $0\le W$ and $L\le W$. Form the
initial and final configurations with states $l$ and $l'$, respectively, phase at its
start, input head at position one, every register tape storing the binary expansion of
its assigned value with head at zero, and output $o$. If the program reaches the final
configuration from the initial one and, at every strictly earlier time, is at an
ordinary (non-call) state with register $r$'s head at some position $q$ satisfying
$-1\le q\le L$, then the run is a simulation to the final configuration in which every
strictly earlier configuration is an ordinary state and every register head position
lies between $-1$ and $W$, inclusive.
\end{lemma}

\begin{lemma}[Simulation of a valid unary–payload input]
\lean{Complexity.LogProg.sim_valP}
\label{Complexity.LogProg.sim_valP}
\leanok
\uses{Complexity.LogProg.HaltsWith, Complexity.LogProg.SimTo, Complexity.LogProg.ValidPlain, Complexity.LogProg.armProg, Complexity.LogProg.armProg_call, Complexity.LogProg.armProg_tr, Complexity.LogProg.aseam, Complexity.LogProg.haltsWith_of_rejects, Complexity.LogProg.simTo_of_reaches, Complexity.LogProg.valPlain_run, Complexity.LogProg.xCfg_aseam, Complexity.LogProg.xCfg_aseam_self}
\difficulty{3}
Let $A$ be a register-program instruction map, let $l_0,l,l'$ be program states, let $r$
be a register, and suppose the instruction at $l$ is the validation instruction for
register $r$ with successor $l'$. For any oracle, register values $v$, output word, and
integer $W\ge 0$, if the input is a valid encoding of a unary prefix paired with a
canonical binary payload, the program simulates from the initial configuration with
state $l$ to the corresponding configuration with state $l'$; if the input is not valid,
it halts with the Boolean answer false. Here canonical means that a nonempty payload
ends in $1$, and the empty payload is allowed.
\end{lemma}

\begin{lemma}[Simulation of a valid nested-pair input]
\lean{Complexity.LogProg.sim_valQ}
\label{Complexity.LogProg.sim_valQ}
\leanok
\uses{Complexity.LogProg.HaltsWith, Complexity.LogProg.SimTo, Complexity.LogProg.ValidPair, Complexity.LogProg.armProg, Complexity.LogProg.armProg_call, Complexity.LogProg.armProg_tr, Complexity.LogProg.aseam, Complexity.LogProg.haltsWith_of_rejects, Complexity.LogProg.simTo_of_reaches, Complexity.LogProg.valPair_run, Complexity.LogProg.xCfg_aseam, Complexity.LogProg.xCfg_aseam_self}
\difficulty{3}
Let $A$ be a register-program instruction map, let $l_0,l,l'$ be program states, let $r$
be a register, and suppose the instruction at $l$ validates a nested pair using register
$r$, with successor $l'$. For any oracle, register values $v$, output word, and integer
$W\ge 0$, if the input encodes a unary prefix paired with two nested payloads, each
canonical (every nonempty payload ends in $1$), the program simulates from the initial
configuration with state $l$ to the corresponding configuration with state $l'$;
otherwise it halts with answer false.
\end{lemma}

\begin{lemma}[Simulation of an input-word equality test]
\lean{Complexity.LogProg.sim_jeqIn}
\label{Complexity.LogProg.sim_jeqIn}
\leanok
\uses{Complexity.LogProg.Ph, Complexity.LogProg.SimTo, Complexity.LogProg.ValidPlain, Complexity.LogProg.armProg, Complexity.LogProg.armProg_call, Complexity.LogProg.armProg_tr, Complexity.LogProg.aseam, Complexity.LogProg.jeqPlain_run, Complexity.LogProg.plainWord, Complexity.LogProg.plainWord_pairEncode, Complexity.LogProg.simTo_of_reachesB, Complexity.LogProg.xCfg_aseam, Complexity.LogProg.xCfg_aseam_self}
\difficulty{4}
Let $A$ be a map from labels to register-program instructions, with initial label $l_0$,
and let an oracle assign a Boolean answer to each decider and bit string. Suppose $x$ is
a bit string of the form $\langle 1^n,w\rangle$, where $w$ is canonical, and $A(l)$
tests whether register $r$ equals the input payload, with successor labels $l_1$ for
equality and $l_0'$ otherwise. For any register valuation $v$, output word, and integer
$W$ such that the length of the binary expansion of $v(r)$ is at most $W$, the compiled
program simulates this instruction from the configuration with label $l$, valuation $v$,
and the given output to the configuration with the same valuation and output, and label
$l_1$ precisely when the binary expansion of $v(r)$ equals $w$ (and $l_0'$ otherwise);
at every step before the endpoint, no oracle call is pending and every work-tape head
position lies between $-1$ and $W$, inclusive.
\end{lemma}

\begin{lemma}[Simulation of a second-segment equality test]
\lean{Complexity.LogProg.sim_jeqSnd}
\label{Complexity.LogProg.sim_jeqSnd}
\leanok
\uses{Complexity.LogProg.Ph, Complexity.LogProg.SimTo, Complexity.LogProg.ValidPair, Complexity.LogProg.armProg, Complexity.LogProg.armProg_call, Complexity.LogProg.armProg_tr, Complexity.LogProg.aseam, Complexity.LogProg.jeqPairSnd_run, Complexity.LogProg.pairWords, Complexity.LogProg.pairWords_pairEncode, Complexity.LogProg.simTo_of_reachesB, Complexity.LogProg.xCfg_aseam, Complexity.LogProg.xCfg_aseam_self}
\difficulty{4}
Let $m$ be the number of registers, let $\Lambda$ be a set of program states, and let
$A$ assign an instruction to each state. Fix an initial state $l_*$, an oracle on
Boolean words, an input word $x$, register values
$v\colon\{0,\ldots,m-1\}\to\mathbb{N}$, an output word, and an integer $W$. Suppose the
instruction at $l$ tests whether register $r$ equals the second component of the pair
decoded from $x$, with successor states $l_1$ and $l_0'$, and suppose the binary
representation of $v(r)$ has length at most $W$ and $x$ is a valid pair input. Then the
compiled program, starting in state $l$ with those register values and output, reaches
the configuration with state $l_1$ if the binary representation of $v(r)$ equals that
decoded second component, and state $l_0'$ otherwise; throughout the run before that
configuration, no call state is visited and every register head position lies between
$-1$ and $W$, inclusive.
\end{lemma}

\begin{lemma}[Simulation of First-Component Equality]
\lean{Complexity.LogProg.sim_jeqFst}
\label{Complexity.LogProg.sim_jeqFst}
\leanok
\uses{Complexity.LogProg.Ph, Complexity.LogProg.SimTo, Complexity.LogProg.ValidPair, Complexity.LogProg.armProg, Complexity.LogProg.armProg_call, Complexity.LogProg.armProg_tr, Complexity.LogProg.aseam, Complexity.LogProg.jeqPairFst_run, Complexity.LogProg.pairWords, Complexity.LogProg.pairWords_pairEncode, Complexity.LogProg.simTo_of_reachesB, Complexity.LogProg.xCfg_aseam, Complexity.LogProg.xCfg_aseam_self}
\difficulty{4}
Let $m,d$ be natural numbers, let $\Lambda$ be a type of labels, and let $A$ assign an
instruction to each label, with initial label $l_0$. Let an oracle assign a Boolean
answer to each decider in a set of $d$ deciders and each bit string, and let $x$ be a
valid nested pair encoding. Suppose $l,l_1,l_0'\in\Lambda$, $r$ is one of the $m$
registers, and the instruction at $l$ tests whether the binary expansion of register $r$
equals the first component of the inner pair, branching to $l_1$ on equality and to
$l_0'$ otherwise. For any register valuation $v$, output word, and integer bound $W$
such that the length of the binary expansion of $v(r)$ is at most $W$, the compiled
program simulates this instruction from the configuration with input $x$, label $l$,
valuation $v$, and the given output word to the configuration with the same input,
valuation, and output, and with label $l_1$ exactly when the binary expansion of $v(r)$
equals the first inner component extracted from $x$, and label $l_0'$ otherwise. At
every intermediate configuration before the endpoint, no oracle call is pending and
every work-tape head position lies between $-1$ and $W$, inclusive.
\end{lemma}
